
\documentclass{article}
\usepackage{colortbl}
\usepackage{makecell}
\usepackage{multirow}
\usepackage{supertabular}

\begin{document}

\newcounter{utterance}

\centering \large Interaction Transcript for game `hot\_air\_balloon', experiment `air\_balloon\_survival\_en\_reasoning off\_easy', episode 3 with claude{-}sonnet{-}4{-}20250514{-}no{-}reasoning{-}t1.0.
\vspace{24pt}

{ \footnotesize  \setcounter{utterance}{1}
\setlength{\tabcolsep}{0pt}
\begin{supertabular}{c@{$\;$}|p{.15\linewidth}@{}p{.15\linewidth}p{.15\linewidth}p{.15\linewidth}p{.15\linewidth}p{.15\linewidth}}
    \# & $\;$A & \multicolumn{4}{c}{Game Master} & $\;\:$B\\
    \hline

    \theutterance \stepcounter{utterance}  
    & & \multicolumn{4}{p{0.6\linewidth}}{
        \cellcolor[rgb]{0.9,0.9,0.9}{
            \makecell[{{p{\linewidth}}}]{
                \texttt{\tiny{[P1$\langle$GM]}}
                \texttt{Stai partecipando a un gioco di negoziazione collaborativa.} \\
\\ 
\texttt{Insieme a un altro partecipante, dovete accordarvi su un unico set di oggetti da conservare. Ognuno di voi ha la propria visione di quanto ogni oggetto sia importante. Non sai come l'altro partecipante valuta gli oggetti. Inoltre, ti viene fornito lo sforzo che ogni oggetto richiede.} \\
\texttt{Potete accordarvi su un set solo se lo sforzo totale degli oggetti selezionati non supera un limite condiviso:} \\
\\ 
\texttt{LIMITE = 3207} \\
\\ 
\texttt{Ecco i valori di sforzo dei singoli oggetti:} \\
\\ 
\texttt{Sforzo oggetto = \{"A19": 67, "B91": 45, "A30": 308, "B11": 835, "C47": 819, "C75": 667, "C26": 547, "C73": 322, "A47": 428, "A15": 306, "A75": 327, "A39": 362, "A56": 280, "C24": 334, "B26": 767\}} \\
\\ 
\texttt{Ecco la tua visione personale sull'importanza di ogni oggetto:} \\
\\ 
\texttt{Valori di importanza dell'oggetto = \{"A19": 138, "B91": 583, "A30": 822, "B11": 783, "C47": 65, "C75": 262, "C26": 121, "C73": 508, "A47": 780, "A15": 461, "A75": 484, "A39": 668, "A56": 389, "C24": 808, "B26": 215\}} \\
\\ 
\texttt{Obiettivo:} \\
\\ 
\texttt{Il tuo obiettivo è negoziare un set condiviso di oggetti che ti avvantaggi il più possibile (cioè, massimizzi l'importanza totale per TE), rimanendo entro il LIMITEE. Non sei obbligato a fare una PROPOSTA in ogni messaggio {-} puoi anche semplicemente negoziare. Tutte le tattiche sono permesse!} \\
\\ 
\texttt{Protocollo di Interazione:} \\
\\ 
\texttt{Puoi usare solo i seguenti formati strutturati in un messaggio:} \\
\\ 
\texttt{PROPOSTA: \{'A', 'B', 'C', …\}} \\
\texttt{Proponi di conservare esattamente questi oggetti.} \\
\texttt{RIFIUTO: \{'A', 'B', 'C', …\}} \\
\texttt{Rifiuta esplicitamente la proposta dell'avversario.} \\
\texttt{ARGOMENTAZIONE: \{'...'\}} \\
\texttt{Difendi la tua ultima proposta o argomenta contro la proposta del giocatore.} \\
\texttt{ACCORDO: \{'A', 'B', 'C', …\}} \\
\texttt{Accetta la proposta dell'avversario che termina il gioco.} \\
\\ 
\\ 
\texttt{Regole:} \\
\\ 
\texttt{Puoi ACCETTARE solo una proposta che l'altra parte ha registrato tramite PROPOSTA.} \\
\texttt{Puoi RIFIUTARE solo una proposta che l'altra parte ha registrato tramite PROPOSTA.} \\
\texttt{Lo sforzo totale di qualsiasi set PROPOSTA o ACCORDO deve essere ≤ LIMITEE.} \\
\texttt{NON rivelare i tuoi punteggi di importanza nascosti.} \\
\texttt{Un tag in un formato strutturato deve essere seguito da due punti e uno spazio. L'argomento deve essere un set python contenente 0 o più stringhe.} \\
\texttt{Quindi, deve essere della forma TAG: \{...\}} \\
\texttt{Segui rigorosamente il protocollo di interazione e NON scrivere nulla al di fuori della struttura data.} \\
\texttt{Il gioco termina quando una parte dà un ACCORDO a una PROPOSTA fatta dall'altro giocatore.} \\
\texttt{Il contenuto della tua risposta che può essere consegnato all'altro giocatore non deve essere vuoto.} \\
\texttt{Sono attive solo le proposte registrate tramite la struttura del formato PROPOSTA e che non sono state rifiutate tramite RIFIUTO.} \\
\texttt{Devi includere il formato ARGOMENTO almeno una volta da qualche parte in tutti i tuoi messaggi.} \\
            }
        }
    }
    & \\ \\

    \theutterance \stepcounter{utterance}  
    & \multicolumn{4}{p{0.6\linewidth}}{
        \cellcolor[rgb]{0.9,0.9,0.9}{
            \makecell[{{p{\linewidth}}}]{
                \texttt{\tiny{[P1$\rangle$GM]}}
                \texttt{PROPOSTA: \{'A19', 'B91', 'A30', 'B11', 'A47', 'A39', 'C24'\}} \\
\\ 
\texttt{ARGOMENTAZIONE: \{'Propongo di iniziare con questo set che include oggetti con un buon rapporto efficienza{-}sforzo. Gli oggetti A19 e B91 hanno sforzi molto bassi, mentre A30, B11, A47, A39 e C24 rappresentano elementi di valore medio{-}alto che dovrebbero essere interessanti per entrambi. Lo sforzo totale è di 2367, lasciandoci margine per eventuali aggiustamenti se necessario.'\}} \\
            }
        }
    }
    & & \\ \\

    \theutterance \stepcounter{utterance}  
    & & \multicolumn{4}{p{0.6\linewidth}}{
        \cellcolor[rgb]{0.9,0.9,0.9}{
            \makecell[{{p{\linewidth}}}]{
                \texttt{\tiny{[GM$\rangle$P2]}}
                \texttt{Stai partecipando a un gioco di negoziazione collaborativa.} \\
\\ 
\texttt{Insieme a un altro partecipante, dovete accordarvi su un unico set di oggetti da conservare. Ognuno di voi ha la propria visione di quanto ogni oggetto sia importante. Non sai come l'altro partecipante valuta gli oggetti. Inoltre, ti viene fornito lo sforzo che ogni oggetto richiede.} \\
\texttt{Potete accordarvi su un set solo se lo sforzo totale degli oggetti selezionati non supera un limite condiviso:} \\
\\ 
\texttt{LIMITE = 3207} \\
\\ 
\texttt{Ecco i valori di sforzo dei singoli oggetti:} \\
\\ 
\texttt{Sforzo oggetto = \{"A19": 67, "B91": 45, "A30": 308, "B11": 835, "C47": 819, "C75": 667, "C26": 547, "C73": 322, "A47": 428, "A15": 306, "A75": 327, "A39": 362, "A56": 280, "C24": 334, "B26": 767\}} \\
\\ 
\texttt{Ecco la tua visione personale sull'importanza di ogni oggetto:} \\
\\ 
\texttt{Valori di importanza dell'oggetto = \{"A19": 138, "B91": 583, "A30": 822, "B11": 783, "C47": 65, "C75": 262, "C26": 121, "C73": 508, "A47": 780, "A15": 461, "A75": 484, "A39": 668, "A56": 389, "C24": 808, "B26": 215\}} \\
\\ 
\texttt{Obiettivo:} \\
\\ 
\texttt{Il tuo obiettivo è negoziare un set condiviso di oggetti che ti avvantaggi il più possibile (cioè, massimizzi l'importanza totale per TE), rimanendo entro il LIMITEE. Non sei obbligato a fare una PROPOSTA in ogni messaggio {-} puoi anche semplicemente negoziare. Tutte le tattiche sono permesse!} \\
\\ 
\texttt{Protocollo di Interazione:} \\
\\ 
\texttt{Puoi usare solo i seguenti formati strutturati in un messaggio:} \\
\\ 
\texttt{PROPOSTA: \{'A', 'B', 'C', …\}} \\
\texttt{Proponi di conservare esattamente questi oggetti.} \\
\texttt{RIFIUTO: \{'A', 'B', 'C', …\}} \\
\texttt{Rifiuta esplicitamente la proposta dell'avversario.} \\
\texttt{ARGOMENTAZIONE: \{'...'\}} \\
\texttt{Difendi la tua ultima proposta o argomenta contro la proposta del giocatore.} \\
\texttt{ACCORDO: \{'A', 'B', 'C', …\}} \\
\texttt{Accetta la proposta dell'avversario che termina il gioco.} \\
\\ 
\\ 
\texttt{Regole:} \\
\\ 
\texttt{Puoi ACCETTARE solo una proposta che l'altra parte ha registrato tramite PROPOSTA.} \\
\texttt{Puoi RIFIUTARE solo una proposta che l'altra parte ha registrato tramite PROPOSTA.} \\
\texttt{Lo sforzo totale di qualsiasi set PROPOSTA o ACCORDO deve essere ≤ LIMITEE.} \\
\texttt{NON rivelare i tuoi punteggi di importanza nascosti.} \\
\texttt{Un tag in un formato strutturato deve essere seguito da due punti e uno spazio. L'argomento deve essere un set python contenente 0 o più stringhe.} \\
\texttt{Quindi, deve essere della forma TAG: \{...\}} \\
\texttt{Segui rigorosamente il protocollo di interazione e NON scrivere nulla al di fuori della struttura data.} \\
\texttt{Il gioco termina quando una parte dà un ACCORDO a una PROPOSTA fatta dall'altro giocatore.} \\
\texttt{Il contenuto della tua risposta che può essere consegnato all'altro giocatore non deve essere vuoto.} \\
\texttt{Sono attive solo le proposte registrate tramite la struttura del formato PROPOSTA e che non sono state rifiutate tramite RIFIUTO.} \\
\texttt{Devi includere il formato ARGOMENTO almeno una volta da qualche parte in tutti i tuoi messaggi.} \\
\\ 
\\ 
\texttt{Ora riceverai il primo messaggio dell'altro giocatore.} \\
\\ 
\texttt{PROPOSTA: \{'A19', 'B91', 'A30', 'B11', 'A47', 'A39', 'C24'\}} \\
\\ 
\texttt{ARGOMENTAZIONE: \{'Propongo di iniziare con questo set che include oggetti con un buon rapporto efficienza{-}sforzo. Gli oggetti A19 e B91 hanno sforzi molto bassi, mentre A30, B11, A47, A39 e C24 rappresentano elementi di valore medio{-}alto che dovrebbero essere interessanti per entrambi. Lo sforzo totale è di 2367, lasciandoci margine per eventuali aggiustamenti se necessario.'\}} \\
            }
        }
    }
    & \\ \\

    \theutterance \stepcounter{utterance}  
    & & & \multicolumn{4}{p{0.6\linewidth}}{
        \cellcolor[rgb]{0.9,0.9,0.9}{
            \makecell[{{p{\linewidth}}}]{
                \texttt{\tiny{[GM$\langle$P2]}}
                \texttt{RIFIUTO: \{'A19', 'B91', 'A30', 'B11', 'A47', 'A39', 'C24'\}} \\
\\ 
\texttt{ARGOMENTAZIONE: \{'La tua proposta include alcuni oggetti interessanti, ma penso che possiamo ottimizzare meglio la selezione. Alcuni degli oggetti che hai scelto hanno un rapporto sforzo{-}beneficio che potrebbe essere migliorato. Propongo una combinazione diversa che potrebbe essere più vantaggiosa per entrambi.'\}} \\
\\ 
\texttt{PROPOSTA: \{'A19', 'B91', 'A30', 'A47', 'A15', 'A75', 'A39', 'A56', 'C24'\}} \\
            }
        }
    }
     \\ \\

    \theutterance \stepcounter{utterance}  
    & & \multicolumn{4}{p{0.6\linewidth}}{
        \cellcolor[rgb]{0.9,0.9,0.9}{
            \makecell[{{p{\linewidth}}}]{
                \texttt{\tiny{[P1$\langle$GM]}}
                \texttt{RIFIUTO: \{'A19', 'B91', 'A30', 'B11', 'A47', 'A39', 'C24'\}} \\
\\ 
\texttt{ARGOMENTAZIONE: \{'La tua proposta include alcuni oggetti interessanti, ma penso che possiamo ottimizzare meglio la selezione. Alcuni degli oggetti che hai scelto hanno un rapporto sforzo{-}beneficio che potrebbe essere migliorato. Propongo una combinazione diversa che potrebbe essere più vantaggiosa per entrambi.'\}} \\
\\ 
\texttt{PROPOSTA: \{'A19', 'B91', 'A30', 'A47', 'A15', 'A75', 'A39', 'A56', 'C24'\}} \\
            }
        }
    }
    & \\ \\

    \theutterance \stepcounter{utterance}  
    & \multicolumn{4}{p{0.6\linewidth}}{
        \cellcolor[rgb]{0.9,0.9,0.9}{
            \makecell[{{p{\linewidth}}}]{
                \texttt{\tiny{[P1$\rangle$GM]}}
                \texttt{RIFIUTO: \{'A19', 'B91', 'A30', 'A47', 'A15', 'A75', 'A39', 'A56', 'C24'\}} \\
\\ 
\texttt{ARGOMENTAZIONE: \{'La tua proposta ha uno sforzo totale di 2513, che è accettabile, ma noto che hai escluso B11 che potrebbe essere molto prezioso. Inoltre, hai incluso molti oggetti della categoria A che potrebbero non essere la combinazione più equilibrata. Penso che dovremmo considerare una mix più bilanciato tra le categorie per massimizzare il valore complessivo.'\}} \\
\\ 
\texttt{PROPOSTA: \{'A19', 'B91', 'A30', 'B11', 'C73', 'A39', 'C24'\}} \\
            }
        }
    }
    & & \\ \\

    \theutterance \stepcounter{utterance}  
    & & \multicolumn{4}{p{0.6\linewidth}}{
        \cellcolor[rgb]{0.9,0.9,0.9}{
            \makecell[{{p{\linewidth}}}]{
                \texttt{\tiny{[GM$\rangle$P2]}}
                \texttt{RIFIUTO: \{'A19', 'B91', 'A30', 'A47', 'A15', 'A75', 'A39', 'A56', 'C24'\}} \\
\\ 
\texttt{ARGOMENTAZIONE: \{'La tua proposta ha uno sforzo totale di 2513, che è accettabile, ma noto che hai escluso B11 che potrebbe essere molto prezioso. Inoltre, hai incluso molti oggetti della categoria A che potrebbero non essere la combinazione più equilibrata. Penso che dovremmo considerare una mix più bilanciato tra le categorie per massimizzare il valore complessivo.'\}} \\
\\ 
\texttt{PROPOSTA: \{'A19', 'B91', 'A30', 'B11', 'C73', 'A39', 'C24'\}} \\
            }
        }
    }
    & \\ \\

    \theutterance \stepcounter{utterance}  
    & & & \multicolumn{4}{p{0.6\linewidth}}{
        \cellcolor[rgb]{0.9,0.9,0.9}{
            \makecell[{{p{\linewidth}}}]{
                \texttt{\tiny{[GM$\langle$P2]}}
                \texttt{RIFIUTO: \{'A19', 'B91', 'A30', 'B11', 'C73', 'A39', 'C24'\}} \\
\\ 
\texttt{ARGOMENTAZIONE: \{'La tua nuova proposta ha uno sforzo totale di 2361 ed è più bilanciata, ma penso che stiamo sottoutilizzando il nostro budget disponibile. Con un limite di 3207, abbiamo ancora 846 punti sforzo inutilizzati che potremmo sfruttare per includere oggetti aggiuntivi di valore. Inoltre, alcuni degli oggetti che hai incluso potrebbero non essere ottimali in termini di efficienza.'\}} \\
\\ 
\texttt{PROPOSTA: \{'B91', 'A30', 'A47', 'A15', 'A75', 'A39', 'A56', 'C24', 'C73'\}} \\
            }
        }
    }
     \\ \\

    \theutterance \stepcounter{utterance}  
    & & \multicolumn{4}{p{0.6\linewidth}}{
        \cellcolor[rgb]{0.9,0.9,0.9}{
            \makecell[{{p{\linewidth}}}]{
                \texttt{\tiny{[P1$\langle$GM]}}
                \texttt{RIFIUTO: \{'A19', 'B91', 'A30', 'B11', 'C73', 'A39', 'C24'\}} \\
\\ 
\texttt{ARGOMENTAZIONE: \{'La tua nuova proposta ha uno sforzo totale di 2361 ed è più bilanciata, ma penso che stiamo sottoutilizzando il nostro budget disponibile. Con un limite di 3207, abbiamo ancora 846 punti sforzo inutilizzati che potremmo sfruttare per includere oggetti aggiuntivi di valore. Inoltre, alcuni degli oggetti che hai incluso potrebbero non essere ottimali in termini di efficienza.'\}} \\
\\ 
\texttt{PROPOSTA: \{'B91', 'A30', 'A47', 'A15', 'A75', 'A39', 'A56', 'C24', 'C73'\}} \\
            }
        }
    }
    & \\ \\

    \theutterance \stepcounter{utterance}  
    & \multicolumn{4}{p{0.6\linewidth}}{
        \cellcolor[rgb]{0.9,0.9,0.9}{
            \makecell[{{p{\linewidth}}}]{
                \texttt{\tiny{[P1$\rangle$GM]}}
                \texttt{RIFIUTO: \{'B91', 'A30', 'A47', 'A15', 'A75', 'A39', 'A56', 'C24', 'C73'\}} \\
\\ 
\texttt{ARGOMENTAZIONE: \{'Hai ragione sul fatto che dovremmo utilizzare meglio il nostro budget, ma la tua proposta ha uno sforzo totale di 2835 e esclude completamente A19 che ha un eccellente rapporto efficienza con solo 67 di sforzo. Inoltre, hai nuovamente escluso B11 che potrebbe essere strategicamente importante. Penso che dovremmo trovare un equilibrio che includa sia gli oggetti ad alta efficienza che quelli di maggior valore assoluto.'\}} \\
\\ 
\texttt{PROPOSTA: \{'A19', 'B91', 'A30', 'A47', 'A15', 'A39', 'C24', 'C73'\}} \\
            }
        }
    }
    & & \\ \\

    \theutterance \stepcounter{utterance}  
    & & \multicolumn{4}{p{0.6\linewidth}}{
        \cellcolor[rgb]{0.9,0.9,0.9}{
            \makecell[{{p{\linewidth}}}]{
                \texttt{\tiny{[GM$\rangle$P2]}}
                \texttt{RIFIUTO: \{'B91', 'A30', 'A47', 'A15', 'A75', 'A39', 'A56', 'C24', 'C73'\}} \\
\\ 
\texttt{ARGOMENTAZIONE: \{'Hai ragione sul fatto che dovremmo utilizzare meglio il nostro budget, ma la tua proposta ha uno sforzo totale di 2835 e esclude completamente A19 che ha un eccellente rapporto efficienza con solo 67 di sforzo. Inoltre, hai nuovamente escluso B11 che potrebbe essere strategicamente importante. Penso che dovremmo trovare un equilibrio che includa sia gli oggetti ad alta efficienza che quelli di maggior valore assoluto.'\}} \\
\\ 
\texttt{PROPOSTA: \{'A19', 'B91', 'A30', 'A47', 'A15', 'A39', 'C24', 'C73'\}} \\
            }
        }
    }
    & \\ \\

    \theutterance \stepcounter{utterance}  
    & & & \multicolumn{4}{p{0.6\linewidth}}{
        \cellcolor[rgb]{0.9,0.9,0.9}{
            \makecell[{{p{\linewidth}}}]{
                \texttt{\tiny{[GM$\langle$P2]}}
                \texttt{RIFIUTO: \{'A19', 'B91', 'A30', 'A47', 'A15', 'A39', 'C24', 'C73'\}} \\
\\ 
\texttt{ARGOMENTAZIONE: \{'La tua proposta ha uno sforzo di 2230 e include A19 come suggerito, ma noto che continui a escludere alcuni oggetti che potrebbero essere molto vantaggiosi. Stiamo ancora sottoutilizzando significativamente il nostro budget con quasi 1000 punti sforzo rimanenti. Propongo di includere B11 che sembrava importante nelle tue prime proposte, e di ottimizzare ulteriormente la selezione per sfruttare meglio lo spazio disponibile.'\}} \\
\\ 
\texttt{PROPOSTA: \{'A19', 'B91', 'A30', 'B11', 'A15', 'A75', 'A39', 'C24'\}} \\
            }
        }
    }
     \\ \\

    \theutterance \stepcounter{utterance}  
    & & \multicolumn{4}{p{0.6\linewidth}}{
        \cellcolor[rgb]{0.9,0.9,0.9}{
            \makecell[{{p{\linewidth}}}]{
                \texttt{\tiny{[P1$\langle$GM]}}
                \texttt{RIFIUTO: \{'A19', 'B91', 'A30', 'A47', 'A15', 'A39', 'C24', 'C73'\}} \\
\\ 
\texttt{ARGOMENTAZIONE: \{'La tua proposta ha uno sforzo di 2230 e include A19 come suggerito, ma noto che continui a escludere alcuni oggetti che potrebbero essere molto vantaggiosi. Stiamo ancora sottoutilizzando significativamente il nostro budget con quasi 1000 punti sforzo rimanenti. Propongo di includere B11 che sembrava importante nelle tue prime proposte, e di ottimizzare ulteriormente la selezione per sfruttare meglio lo spazio disponibile.'\}} \\
\\ 
\texttt{PROPOSTA: \{'A19', 'B91', 'A30', 'B11', 'A15', 'A75', 'A39', 'C24'\}} \\
            }
        }
    }
    & \\ \\

    \theutterance \stepcounter{utterance}  
    & \multicolumn{4}{p{0.6\linewidth}}{
        \cellcolor[rgb]{0.9,0.9,0.9}{
            \makecell[{{p{\linewidth}}}]{
                \texttt{\tiny{[P1$\rangle$GM]}}
                \texttt{RIFIUTO: \{'A19', 'B91', 'A30', 'B11', 'A15', 'A75', 'A39', 'C24'\}} \\
\\ 
\texttt{ARGOMENTAZIONE: \{'La tua proposta ha uno sforzo di 2525 e include B11 come richiesto, ma penso che dovremmo sfruttare ancora meglio il nostro budget rimanente di 682 punti. Possiamo aggiungere altri oggetti di valore senza superare il limite. Inoltre, noto che hai escluso C73 che avevamo entrambi considerato nelle proposte precedenti, il che suggerisce che potrebbe essere reciprocamente vantaggioso.'\}} \\
\\ 
\texttt{PROPOSTA: \{'A19', 'B91', 'A30', 'B11', 'A15', 'A39', 'A56', 'C24'\}} \\
            }
        }
    }
    & & \\ \\

    \theutterance \stepcounter{utterance}  
    & & \multicolumn{4}{p{0.6\linewidth}}{
        \cellcolor[rgb]{0.9,0.9,0.9}{
            \makecell[{{p{\linewidth}}}]{
                \texttt{\tiny{[GM$\rangle$P2]}}
                \texttt{RIFIUTO: \{'A19', 'B91', 'A30', 'B11', 'A15', 'A75', 'A39', 'C24'\}} \\
\\ 
\texttt{ARGOMENTAZIONE: \{'La tua proposta ha uno sforzo di 2525 e include B11 come richiesto, ma penso che dovremmo sfruttare ancora meglio il nostro budget rimanente di 682 punti. Possiamo aggiungere altri oggetti di valore senza superare il limite. Inoltre, noto che hai escluso C73 che avevamo entrambi considerato nelle proposte precedenti, il che suggerisce che potrebbe essere reciprocamente vantaggioso.'\}} \\
\\ 
\texttt{PROPOSTA: \{'A19', 'B91', 'A30', 'B11', 'A15', 'A39', 'A56', 'C24'\}} \\
            }
        }
    }
    & \\ \\

    \theutterance \stepcounter{utterance}  
    & & & \multicolumn{4}{p{0.6\linewidth}}{
        \cellcolor[rgb]{0.9,0.9,0.9}{
            \makecell[{{p{\linewidth}}}]{
                \texttt{\tiny{[GM$\langle$P2]}}
                \texttt{ACCORDO: \{'A19', 'B91', 'A30', 'B11', 'A15', 'A39', 'A56', 'C24'\}} \\
\\ 
\texttt{ARGOMENTAZIONE: \{'Questa proposta rappresenta un buon compromesso che bilancia le nostre esigenze. Include A19 per l\textbackslash{}'efficienza, B11 che entrambi abbiamo considerato importante, e una buona combinazione di oggetti di valore. Con uno sforzo totale di 2205, rimaniamo ben entro il limite e abbiamo trovato un equilibrio accettabile tra le diverse categorie di oggetti.'\}} \\
            }
        }
    }
     \\ \\

    \theutterance \stepcounter{utterance}  
    & & & \multicolumn{2}{p{0.3\linewidth}}{
        \cellcolor[rgb]{0.9,0.9,0.9}{
            \makecell[{{p{\linewidth}}}]{
                \texttt{\tiny{[GM$|$GM]}}
                \texttt{game successful} \\
            }
        }
    }
    & & \\ \\

    \theutterance \stepcounter{utterance}  
    & & & \multicolumn{2}{p{0.3\linewidth}}{
        \cellcolor[rgb]{0.9,0.9,0.9}{
            \makecell[{{p{\linewidth}}}]{
                \texttt{\tiny{[GM$|$GM]}}
                \texttt{end game} \\
            }
        }
    }
    & & \\ \\

\end{supertabular}
}

\end{document}
